% ex: ts=2 sw=2 sts=2 et filetype=tex
% SPDX-License-Identifier: CC-BY-SA-4.0
\input{cv/header}
%%%%%%%%%%%%%%%%%%%%%%%%% Begin CV Document %%%%%%%%%%%%%%%%%%%%%%%%%%%%

\begin{document}
\makeheading{Miguel Bernal Marin}

\section{Contacto}
%
% NOTE: Mind where the & separators and \\ breaks are in the following
%       table.
%
% ALSO: \rcollength is the width of the right column of the table
%       (adjust it to your liking; default is 1.85in).
%
\newlength{\rcollength}\setlength{\rcollength}{2.85in}%
%
\begin{tabular}[t]{@{}p{\textwidth-\rcollength}p{\rcollength}}
  \textit{Teléfono:} \IfFileExists{cv/phone.tex}%
  {\input{cv/phone.tex}}
  {\href{https://wa.me/52XXX}{+52 XXX}}
 &
  \textit{E-mail:} \email{miguel.bernal.marin@gmail.com}\\
Telegram: \href{https://t.me/miguelinux}{@miguelinux} &
  \href{https://scholar.google.com/citations?user=-9GezEIAAAAJ\&hl=es}
  {Google Scholar/user=-9GezEIAAAAJ} \\
Zapopan, Jalisco, México  &
	\href{https://github.com/miguelinux}{GitHub.com/miguelinux} \\
   & \href{https://www.linkedin.com/in/miguelinux/}
	{Linkedin.com/in/miguelinux} \\
\end{tabular}
%\section{Nacionalidad}
%
%Mexicana

\section{Intereses Personales}
%
Linux, sysadmin, devops, 
programación de computadoras,
educación,
visión por computadora,
robótica,
robots móviles,
realidad virtual,
diseño de sistemas.


\section{Educación}
%
\href{https://unidad.gdl.cinvestav.mx/}{\textbf{Centro de Investigación y Estudios Avanzados del I.P.N }}(Cinvestav \emph{Unidad Guadalajara}),
Zapopan, Jalisco, México

\begin{outerlist}
\item[] Doctorado en Ciencias en Ingeniería Eléctrica, \hfill{Febrero 2012}
        \begin{innerlist}
        \item Tema de tesis: \emph{Técnicas geométricas y aprendizaje automático para la
              reconstrucción 3D y navegación robótica}
        \item Asesor:
              \href{https://unidad.gdl.cinvestav.mx/investigadores/investigador.php?inv=11}
                   {Prof. Dr. Eduardo Bayro~Corrochano}
        \item Área de Estudio: Visión por computadora \& robots móviles
        \end{innerlist}

\item[] Maestría en Ciencias en Ingeniería Eléctrica,\hfill{Agosto 2007}
        \begin{innerlist}
        \item Tema de tesis: \emph{Mapas tridimencionales para navegación
              robótica}
        \item Asesor:
              \href{https://unidad.gdl.cinvestav.mx/investigadores/investigador.php?inv=11}
                   {Prof. Dr. Eduardo Bayro~Corrochano}
        \item Especialidad: \emph{Ciencias Computacionales}
        \item Área de Estudio: Visión por computadora \& robots móviles
        \end{innerlist}
\end{outerlist}

\blankline

\href{http://www.udg.mx/}{\textbf{Universidad de Guadalajara}},
Guadalajara, Jalisco, México

\begin{outerlist}
\item[] \href{http://www.cucei.udg.mx/es/oferta-academica/licenciaturas/licenciatura-en-matematicas }
	{Licenciatura en Matemáticas},
\hfill{Mayo 2002}
        \begin{innerlist}
        \item Especialización en computación matemática
        \end{innerlist}
\end{outerlist}

\section{Estancias Extranjero}
%
\textbf{Fraunhofer IITB},
Karlsruhe, Alemania
\begin{outerlist}
\item[] \textit{Estancia de Investigación}%
    \hfill \textbf{Mayo 2008 a Noviembre 2008}
    \begin{innerlist}
        \item Tema: Uso de álgebras geométricas y aprendizaje automático
              en un robot móvil
        \item Supervisor: Dr.-Ing. Helge-Bj{\"o}rn Kuntze
    \end{innerlist}
\end{outerlist}


\section{Docencia}
% ex: ts=2 sw=2 sts=2 et filetype=tex
% SPDX-License-Identifier: CC-BY-SA-4.0

\href{https://www.tecmm.edu.mx/}
{\textbf{Instituto Tecnológico Superior de Jalisco}},
Zapopan, Jalisco, México
\begin{outerlist}
\item[] \textit{Profesor de Maestría}%
    \hfill \textbf{Sep 2012 - Actual }
    \begin{innerlist}
        \item Gráficas por computadora, Análisis y Diseño de Algoritmos,
	      Imágenes Digitales, Sistemas Distribuidos.
    \end{innerlist}
\item[] \textit{Profesor}%
    \hfill \textbf{Mar/07 a Feb/08, Mar/09 - Actual}
    \begin{innerlist}
        \item Graficación, Sistemas Operativos, Estructura de Datos,
		Programación Orientada a Objecto, Fundamentos de Programación,
		Tópicos de Seguridad Web, Base de datos, Diseño y Análisis
		de Algoritmos, Ciberseguridad.
    \end{innerlist}
\end{outerlist}

\blankline

\href{https://iteso.mx/}
{\textbf{Instituto Tecnológico y de Estudios Superiores de Occidente (ITESO)}},
Tlaquepaque, Jalisco, México
\begin{outerlist}
\item[] \textit{Profesor de asignatura}%
    \hfill \textbf{Primavera 2020 - Actual}
    \begin{innerlist}
        \item Algoritmos y Programación, Seguridad en Redes
    \end{innerlist}
\end{outerlist}

\blankline

\href{https://tec.mx/}
{\textbf{Tecnológico de Monterrey}},
Zapopan, Jalisco, México
\begin{outerlist}
\item[] \textit{Teacher}%
    \hfill \textbf{Sep 2021 - Actual}
    \begin{innerlist}
        \item Computo en la nube, visión computacional
    \end{innerlist}
\end{outerlist}

\blankline

\href{http://www.colomos.ceti.mx/}{\textbf{Centro de Enseñanza Técnica
Industrial}} (CETI plantel Colomos),
Guadalajara, Jalisco, México
\begin{outerlist}
\item[] \textit{Profesor}%
    \hfill \textbf{Marzo 2007 a Febrero 2008}
    \begin{innerlist}
        \item Computación IV
    \end{innerlist}
\end{outerlist}

\blankline

\href{https://www.ucg.edu.mx/}{\textbf{Universidad Cuauhtémoc} plantel Guadalajara},
Zapopan, Jalisco, México
\begin{outerlist}
\item[] \textit{Profesor}%
    \hfill \textbf{Mayo 2005 a Diciembre 2007}
    \begin{innerlist}
        \item Sistemas Operativos, Linux, Supercomputo, Solaris, Taller de
Programación Java, Programación para Internet.
    \end{innerlist}
\item[] \textit{Profesor Sistema Empresarial}%
    \hfill \textbf{Enero 2006 a Abril 2008}
    \begin{innerlist}
        \item Sistemas Operativos, Linux, Supercomputo, Solaris, Taller de
Programación Java, Programación para Internet, Métodos numéricos, Cálculo
Integral, Cálculo Diferencial, Matemáticas I, Matemáticas II.
    \end{innerlist}
\item[] \textit{Profesor Maestría en Sistemas Computacionales}%
    \hfill \textbf{Febrero 2006}
    \begin{innerlist}
        \item Paradigmas de programación
    \end{innerlist}
\end{outerlist}

\blankline

\href{http://www.udg.mx/}{\textbf{Universidad de Guadalajara}},
Guadalajara, Jalisco, México
\begin{outerlist}
\item[] \textit{Profesor suplente}%
    \hfill \textbf{Febrero 2002 a Agosto 2002}
    \begin{innerlist}
        \item Geometría Analítica
    \end{innerlist}
\end{outerlist}



\blankline

\section{Patentes}
% ex: ts=2 sw=2 sts=2 et filetype=tex
% SPDX-License-Identifier: CC-BY-SA-4.0

\begin{bibsection}

  \item Jose Camacho Perez, Miguel Bernal Marin, Mario Alfredo Carrillo
	Arevalo, Hector Cordourier Maruri, Jesus Adan Cruz Vargas, Abraham
	Duenas De La Cruz, Paulo Lopez Mayer, Julio Zamora Esquivel.
	Surface detection for mobile devices.
	\emph{US Patent 10,609,205 B2}, Mar. 31, 2020

  \item Jose Camacho Perez, Miguel Bernal Marin, Mario Alfredo Carrillo
	Arevalo, Hector Cordourier Maruri, Jesus Adan Cruz Vargas, Abraham
	Duenas De La Cruz, Paulo Lopez Mayer, Julio Zamora Esquivel.
	Surface detection for mobile devices.
	\emph{US Patent 10,938,977 B2}, Mar. 2, 2021

  \item Jose Camacho Perez, Miguel Bernal Marin, Mario Alfredo Carrillo
	Arevalo, Hector Cordourier Maruri, Jesus Adan Cruz Vargas, Abraham
	Duenas De La Cruz, Paulo Lopez Mayer, Julio Zamora Esquivel.
	Surface detection for mobile devices.
	\emph{US Patent 11,516,336 B2}, Nov. 9, 2022

\end{bibsection}

\blankline

\section{Publicación Revistas}
% ex: ts=2 sw=2 sts=2 et filetype=tex
% SPDX-License-Identifier: CC-BY-SA-4.0

\begin{bibsection}
  \item Alberto Petrilli-Barceló, Heriberto Casarrubias-Vargas,
        Miguel Bernal-Marin, Eduardo Bayro-Corrochano, R{\"u}diger Dillmann.
        Geometric Techniques for Humanoid Perception,
        \emph{International Journal of Humaoid Robotics},
        vol. 7, no 3, pp.429-450, 2010.

  \item Miguel Bernal-Marin, Eduardo Bayro-Corrochano.
        Integration of Hough Transform of lines and planes in the
        framework of conformal geometric algebra for 2D and 3D robot vision.
        \emph{Pattern Recognition Letters},
        vol. 32, no 16, pp. 2213-2223, 2011.
\end{bibsection}


\blankline

\section{Capítulos de libros}
% ex: ts=2 sw=2 sts=2 et filetype=tex
% SPDX-License-Identifier: CC-BY-SA-4.0

\begin{bibsection}
  \item Miguel Bernal-Marin, Adan Landa-Hernández, Eduardo Bayro-Corrochano,
        Chapter 7 - Visual Geometric Entities for Orientation,
        Relocalization and Navigation,
        in \emph{Autonomous Robots: Control, Sensing and Perception}.
        pp. 152-183, A. Torres, E. Martínez, Eds. Cuvillier Verlag, 2011.
        ISBN 978-3-86955-866-0.
\end{bibsection}


\section{Publicación Conferencias}
% ex: ts=2 sw=2 sts=2 et filetype=tex
% SPDX-License-Identifier: CC-BY-SA-4.0

\begin{bibsection}
  \item Eduardo Bayro-Corrochano, Miguel Bernal-Marin.
        Generalized Hough Transform and Conformal Geometric Algebra
        to Detect Lines and Planes for  Building 3D Maps and Robot Navigation
        \emph{IEEE/RSJ International Conference on Intelligent RObots and
        Systems}, Octubre, 2010.

  \item Miguel Bernal-Marin, Eduardo Bayro-Corrochano.
        Geometric Hough transform for robot vision,
        \emph{Applied Geometric Algebras in Computer Science and Engineering},
        Junio 2010.

  \item Miguel Bernal-Marin, Eduardo Bayro-Corrochano.
        Machine learning and geometric technique for SLAM,
        \emph{Iberoamerican Congress on Pattern Recognition}, Noviembre, 2009.

  \item Miguel Bernal-Marin, Eduardo Bayro-Corrochano.
        Visual and Laser Guided Robot Relocalization Using Lines,
        Hough Transformation and Machine Learning Techniques,
        \emph{IEEE/RSJ International Conference on Intelligent RObots and
        Systems}, Octubre, 2009.

  \item Miguel Bernal-Marin, Eduardo Bayro-Corrochano.
        Visual and Laser Guided Robot Relocalization Using Lines and Hough
        Transformation, \emph{IEEE-RAS International Conference on
        Humanoid Robots}, Diciembre, 2008.
\end{bibsection}


\section{Publicaciones Internas de Intel}
% ex: ts=2 sw=2 sts=2 et filetype=tex
% SPDX-License-Identifier: CC-BY-SA-4.0

\begin{bibsection}
  \item Julio Zamora-Esquivel, Miguel Bernal-Marin, Omar Estrada-Diaz.
	Boosting 60X presillicon-setup for execution,
	\emph{Software Professionals Conference} (SWPC), 2013

  \item Julio Zamora-Esquivel, Miguel Bernal-Marin, Alberto Castrejon.
        Unveiling PixWin image compare algorithm,
        \emph{VPG Tech Submmit (VTS)}, 2012.

  \item Julio Zamora-Esquivel, Miguel Bernal-Marin, Alberto Castrejon.
        Best Presentation Demo: \emph{An automation tool called Helix},
        \emph{Sofware Professionals Conference (SWPC)}, 2012.
\end{bibsection}

\blankline

\section{Experiencia Profesional}
%
Intel Tecnología de México,
%\textsuperscript{\textregistered}
Zapopan, Jalisco, México
\begin{outerlist}
\item[] \textit{Ingeniero de desarrollo de sistemas operativos}%
    \hfill \textbf{Agosto 2021 - Actual}
    \begin{innerlist}
      \item Habilitador de las características de los procesadores Intel en
	      la distribución de Linux Debian
      \item Portar herramientas internas y externas a Debian Linux
      \item Desarrollar herramientas en Python para la validación de
        distribuciones de Linux
    \end{innerlist}
\item[] \textit{Ingeniero de Software}%
    \hfill \textbf{Agosto 2014 - Agosto 2021}
    \begin{innerlist}
      \item Desarrollador de Linux en el proyecto de
	      \href{https://www.clearlinux.org/}{Clear Linux}
      \item Mantenedor del paquete del kernel de Linux en Clear Linux
    \end{innerlist}
\item[] \textit{Ingeniero de Gráficos por Software}%
    \hfill \textbf{Agosto 2011 a Agosto 2014}
    \begin{innerlist}
        \item Desarrollador y Validador de pruebas del API de
           Microsoft\textsuperscript{\textregistered}
           DirectX\textsuperscript{\textregistered}.
        \item Validador de gráficos en Pre-silicio usando Fulsim, HAS 2.0 \& 3.0
	   sobre una arquitectura Linux Xen
	\item Validador de gráficos en Android en arquitectura Intel
    \end{innerlist}
\end{outerlist}

\blankline

Promotora Atletica Speedy, Guadalajara, Jalisco, México
\begin{outerlist}
\item[] \textit{Webmaster}%
    \hfill \textbf{Febrero 2003 - Agosto 2004 }
    \begin{innerlist}
        \item Desarrollo del sitio http://www.promospeedy.com
    \end{innerlist}
\end{outerlist}

\blankline

Sistemas Integrales a la Medida, Zapopan, Jalisco, México
\begin{outerlist}
\item[] \textit{Webmaster }%
    \hfill \textbf{Septiembre 1999 - Diciembre 2002 }
    \begin{innerlist}
        \item Desarrollador de sitios web
              (http://www.radaenterprises.com) \\
              (http://www.tocomarketing.com)
        %\item Desarrollador de Base de Datos
        %\item Desarrollo del sitio de la empresa Rada Enterprice
              %(http://www.radaenterprises.com)
        %\item Desarrollo del sitio de la empresa Toco Marketing
              %(http://www.tocomarketing.com)
    \end{innerlist}
\end{outerlist}

\blankline

\section{Expositor}
% ex: ts=2 sw=2 sts=2 et filetype=tex
% SPDX-License-Identifier: CC-BY-SA-4.0

\begin{loneinnerlist}
  %\item  \emph{} \hfill{\emph{2023}}
  \item Introducción a los Quadlet de Podman,
	  \emph{PosaDev} \hfill{\emph{2023}}
  \item Introducción a los contenedores de Linux, \emph{UdG} \hfill{\emph{2023}}
  \item Sistemas de prototipo de chatbot de inteligencia artificial,
	  \emph{SSJ} \hfill{\emph{2023}}
  \item De que están hechas las nubes, \emph{CECITEJ} \hfill{\emph{2023}}
  \item Autenticación a 2 pasos: Qué es y para que sirve, \emph{GUGLAG}
    \hfill{\emph{2022}}
  \item Papà, ¿De que están hechas las nubes?, \emph{UdG} \hfill{\emph{2022}}
  \item Como contribuir al kernel de Linux, \emph{CCOSS} \hfill{\emph{2022}}
  \item Linux en la vida cotidiana, \emph{PosaDev} \hfill{\emph{2021}}
  \item Aplicaciones de la Suite de Google, \emph{Instructor Tec MM} \hfill{\emph{2020}}
  \item G Suite for Education, \emph{Instructor Tec MM} \hfill{\emph{2020}}
  \item Clear Linux, la distribución de Intel, \emph{TopTamaulipas 2019 (Cinvestav Tamaulipas)} \hfill{\emph{2019}}
  \item Seguridad para personas normales, \emph{BSides Guadalajara} \hfill{\emph{2019}}
  \item Taller de Clear Linux, \emph{Cumbre de Contribuidores del Open Source Software} \hfill{\emph{2019}}
  \item ¿Que es OpenCV y para que sirve?, \emph{Festival de Software Libre} \hfill{\emph{2019}}
  \item Clear Linux, la distribución de Intel, \emph{Festival de Software Libre} \hfill{\emph{2018}}
  \item Juez en los concursos de Robots Clasificadores y Mini Sumo,
        \emph{Tecnológico de Monterrey} \hfill{\emph{2010}}
  \item Visión robótica: Robots móviles y humanoides, \emph{Universidad del
        Valle de México} \hfill{\emph{2010}}
  \item Visión Computacional: Cómo perciben el mundo los robots,
        \emph{Instituto Tecnológico de la Laguna} \hfill{\emph{2010}}
  \item Robots Moviles y Humanoides, \emph{Universitronica 2010}
        \hfill{\emph{2010}}
  \item \emph{IEEE/RSJ International Conference on Intelligent RObots
        and Systems}, \hfill{\emph{2009}}
  \item Linux como estación de trabajo y linux como servidor,
        CUValles, \emph{Universidad de Guadalajara} \hfill{\emph{2008}}
  \item Comité Organizador, \emph{Primer Olimpiada Interna de Matemáticas},
        Universidad Cuauh\-témoc plantel Guadalajara \hfill{\emph{2006}}
  \item Comité Organizador Local, \emph{XXXII Congreso Nacional de la  Sociedad
        Matemática Mexicana} \hfill{\emph{1999}}
\end{loneinnerlist}


\section{Contribuciones al open source}
% ex: ts=2 sw=2 sts=2 et filetype=tex
% SPDX-License-Identifier: CC-BY-SA-4.0
%Contribuidor de parches para proyectos de software de código abierto (Open
%Source).
\begin{loneinnerlist}
    \item \href{https://git.kernel.org/pub/scm/linux/kernel/git/torvalds/linux.git/log/?qt=author&q=Miguel+Bernal+Marin}
	    {Linux kernel}
    \item \href{https://salsa.debian.org/kernel-team/linux/-/commits/master?author=Miguel\%20Bernal\%20Marin}
	    {Debian Linux kernel}
    \item \href{https://ivt.sourceforge.net/about.html}{IVT} (The Integrating
      Vision Toolkit)
    %\item \href{http://opencv.willowgarage.com/wiki/}{OpenCV} (Open Source Computer Vision)
    \item \href{https://github.com/pmem/ndctl/commits/?author=miguelinux}
      {ndctl} (non-volatile memory decive control)
    \item \href{https://github.com/systemd/systemd/commits?author=miguelinux}{systemd}
	    (System and Service Manager)
     \item \href{https://github.com/khenidak/dysk/commits/?author=miguelinux}
       {dysk} (Attach Azure disks)
    \item \href{https://github.com/intel/qpl/commits/?author=miguelinux}
      {Intel Query Processing Library }
    \item \href{https://github.com/clearlinux-pkgs/linux/commits/?author=miguelinux}
      {Clear Linux kernel}
\end{loneinnerlist}
%Revisor de propuestas técnicas para la conferencia de profecionales del Software
%de Intel 2013, 2015


\section{Capacitación}
% ex: ts=2 sw=2 sts=2 et filetype=tex
% SPDX-License-Identifier: CC-BY-SA-4.0

\begin{loneinnerlist}
  %\item  \emph{} \hfill{\emph{}}
  \item Diseño de asignaturas en línea, \emph{ITESO} \hfill{\emph{2020}}
  \item Aulas virtuales y Herramientas multimedia, \emph{Tec MM} \hfill{\emph{2020}}
  \item Android Internals, \emph{Marakana - Intel} \hfill{\emph{2012}}
  \item Linux Kernel Internals, \emph{Linux Fundation - Intel} \hfill{\emph{2012}}
  \item Linux Devices Driver, \emph{Start Journey - Intel} \hfill{\emph{2011}}
  \item Educación centrada en el aprendizaje (Diplomado DOCA),
        \emph{Instituto Tecnológico Superior de Zapopan} \hfill{\emph{2009}}
  \item Python, \emph{Universidad de Guadalajara}\hfill{\emph{2003}}
  \item Behaviors on IPV6 for integration with IPV4, \emph{International
        symposium on Advanced Distributed System} \hfill{\emph{2002}}
  \item 2\textsuperscript{do} semestre, \emph{Cisco Networking Academy
        Program} \hfill{\emph{2002}}
  \item 1\textsuperscript{er} semestre, \emph{Cisco Networking Academy
        Program} \hfill{\emph{2001}}
  \item Software Didáctico: "Desarrollo de Software Educativo, una Labor
        eminentemente Interdisciplinaria", \emph{Centro de Investigación
        en Matemáticas, A.C.} (CIMAT) \hfill{\emph{2001}}
  \item Procesamiento de Imágenes, \emph{Centro de Investigación
        en Matemáticas, A.C.} (CIMAT) \hfill{\emph{2001}}
  \item JAVA Básico, \emph{Universidad de Guadalajara}\hfill{\emph{2001}}
  \item Operador de Equipo de Computo, \emph{Instituto de Computación de
        Occidente} \hfill{\emph{1998}}
  \item Capturista de Datos, \emph{Centro de Capacitación de Sistemas
        Contable, Administrativos y de Cómputo, A.C.} \hfill{\emph{1997}}
  \item Internet, \emph{Universidad de Guadalajara} \hfill{\emph{1997}}
\end{loneinnerlist}

\blankline

\section{Concursos congresos} \begin{loneinnerlist}
  %\item \hfill{\emph{}}
  \item \emph{Iberoamerican Congress on Pattern Recognition} \hfill{\emph{2009}}
  \item \emph{IEEE/RSJ International Conference on Intelligent RObots
        and Systems}, \hfill{\emph{2009}}
  \item 1\textsuperscript{er} Festival GNU/Linux y software libre
        \hfill{\emph{2002}}
  \item Días de Software Libre \hfill{\emph{2002}}
  \item \emph{XXXIV} Congreso Nacional de la  Sociedad Matemática
        Mexicana \hfill{\emph{2001}}
  \item \emph{XXXIII} Congreso Nacional de la  Sociedad Matemática
        Mexicana \hfill{\emph{2000}}
  \item \emph{XXXII} Congreso Nacional de la  Sociedad Matemática
        Mexicana \hfill{\emph{1999}}
  \item Concurso por Equipos de Matemáticas \hfill{\emph{1997}}
  \item Olimpiada de Matemáticas \hfill{\emph{1997}}
  \item 2\textsuperscript{do} Concurso Regional de Matemáticas por Equipos
        \hfill{\emph{1997}}
\end{loneinnerlist}

\blankline

\section{Idiomas}
%
Ingles \textit{Nivel Medio Alto}\\%
Aleman \textit{Nivel Básico}%

\section{Habilidades Técnicas}
%
Amplia experiencia en hardware y software de computadoras, redes
informáticas, tecnologías de la información y Linux

\blankline

Programación: C, C$+$$+$, Python, Java, JavaScript, Pascal, PHP, ASP, UNIX
        shell scripting, GNU make, SQL.

\blankline

Tecnologías de la Información: Redes (UDP, TCP, ARP, DNS),
Servicios (HTTP, SQL, POP, IMAP, SMTP)

\blankline

Aplicaciones Computacionales: \TeX{} (\LaTeX{}, B\textsc{ib}\TeX{}),
        aplicaciones comunes de computación (para las plataformas Windows,
        Mac OS X, y Linux), Vim

\blankline

Sistemas Operativos: familia de Microsoft Windows, Apple OS X, Linux, FreeBSD,
        Solaris, y otras variantes de Unix.



\end{document}

%%%%%%%%%%%%%%%%%%%%%%%%%% End CV Document %%%%%%%%%%%%%%%%%%%%%%%%%%%%%
%----------------------------------------------------------------------%
% The following is copyright and licensing information for
% redistribution of this LaTeX source code; it also includes a liability
% statement. If this source code is not being redistributed to others,
% it may be omitted. It has no effect on the function of the above code.
%----------------------------------------------------------------------%
% Copyright (c) 2007, 2008, 2009, 2010, 2011 by Theodore P. Pavlic
%
% Unless otherwise expressly stated, this work is licensed under the
% Creative Commons Attribution-Noncommercial 3.0 United States License. To
% view a copy of this license, visit
% http://creativecommons.org/licenses/by-nc/3.0/us/ or send a letter to
% Creative Commons, 171 Second Street, Suite 300, San Francisco,
% California, 94105, USA.
%
% THE SOFTWARE IS PROVIDED "AS IS", WITHOUT WARRANTY OF ANY KIND, EXPRESS
% OR IMPLIED, INCLUDING BUT NOT LIMITED TO THE WARRANTIES OF
% MERCHANTABILITY, FITNESS FOR A PARTICULAR PURPOSE AND NONINFRINGEMENT.
% IN NO EVENT SHALL THE AUTHORS OR COPYRIGHT HOLDERS BE LIABLE FOR ANY
% CLAIM, DAMAGES OR OTHER LIABILITY, WHETHER IN AN ACTION OF CONTRACT,
% TORT OR OTHERWISE, ARISING FROM, OUT OF OR IN CONNECTION WITH THE
% SOFTWARE OR THE USE OR OTHER DEALINGS IN THE SOFTWARE.
%----------------------------------------------------------------------%

